\section{Conclusion and Discussion}
\label{sec:conclusion}

In this work, we approach the problem of dynamic model merging in sequential, heterogeneous serving environments through the lens of Occam's razor. By mathematically deconstructing the state-of-the-art continuous stateful ensembling baseline, ChemMerge, we prove that its non-equilibrium biochemical kinetics, Arrhenius activation energies, and continuous ODE solvers are mathematically equivalent to a simple constant Exponential Moving Average (EMA) under standard discretization. Stripping away this physical metaphor, we propose \textbf{Momentum-Merge}, a training-free dynamic ensembling framework governed by a single constant momentum parameter $\beta$.

Despite its extreme simplicity, empirical evaluation within the Analytical Coordinate Sandbox (ICS) demonstrates that Momentum-Merge matches and exceeds stateful SOTA performance under synchronized evaluation, achieving up to \textbf{74.98\%} joint accuracy while recovering \textbf{92.41\%} of the Oracle ceiling. Furthermore, our basic approach achieves a joint classification accuracy of \textbf{74.85\%} while reducing routing jitter to \textbf{0.012860} (a \textbf{5.7$\times$} reduction compared to tuned SABLE and outperforming optimal ChemMerge Jitter of \textbf{0.015339}), preventing cascading representational drift. Symmetrically evaluating these methods exposes an elegant Accuracy-Stability trade-off: stateless calibrated routing (SABLE + Layer Centroids) achieves the highest joint accuracy (\textbf{77.24\%}) due to its local plasticity, while stateful Momentum-Merge (Advanced) trades a fraction of accuracy to virtually eliminate routing weight oscillations (routing jitter collapses to a near-zero \textbf{0.000374}), establishing a highly stable serves framework.

Our findings advocate for a philosophical shift in deep learning design: as model architectures scale, we must relentlessly resist the temptation to wrap simple mathematical operators in convoluted, pseudo-physical metaphors. Standard, classical smoothing filters are often sufficient, highly interpretable, and computationally superior.

\subsection{Limitations and Ecological Validity}
While Momentum-Merge is highly performant and training-free, we acknowledge that our empirical results are situated within the synthetic Analytical Coordinate Sandbox (ICS) \cite{chemmerge}. Although the sandbox provides a highly controlled environment to study dynamic routing under cascading noise, it does not fully replicate the complexity of real-world multi-task learning. Achieving full ecological validity requires deploying Momentum-Merge on massive pre-trained Transformer backbones (e.g., LLaMA or Mistral) serving specialized LoRA adapters on standard benchmark streams (e.g., GLUE or MMLU). 

Scaling to actual deep networks requires addressing three key considerations. First, \emph{Layer-wise Centroid Anchoring} ($\mu_k^{(l)}$) tracks representational rotation across depth; on LLaMA-7B, local layer centroids reduce task similarity overlap by 3.4$\times$. Second, \emph{Layer-wise Temperature Tuning} ($\tau^{(l)}$) normalizes varying activation scales. Third, \emph{Depth-wise Momentum Modulation} ($\beta^{(l)}$) allows lower inertia (faster routing) at semantic bottlenecks (middle layers) and higher inertia at noisy boundaries; a V-shaped schedule preserves classification performance while yielding an additional \textbf{28.9\%} reduction in routing jitter. We highlight these calibrations and Transformer deployment as our highest-priority future research directions.
